\scanpage{136}
% archive foliation: 135
\begin{equation}
\mathbb Z^2\longrightarrow H
\tag{18}
\end{equation}
compatible avec les orientations, i.e. tel que l'on ait
\begin{equation}
z_1\wedge z_2=\omega
\tag{19}
\end{equation}
où $z_1,z_2$ sont les images par (18) des éléments de la base canonique de $\mathbb Z^2$. Cet isomorphisme (18) définit un isomorphisme correspondant
\begin{equation}
\mathbb R^2\simeq V
\tag{20}
\end{equation}
d'où des isomorphismes
\begin{equation}
\left\{
\begin{aligned}
\mathbb P_{\mathbb R}^1&\simeq\mathbb P^1(V),&
\mathbb P^1(\mathbb C)&\simeq\Sigma_V\\
\mathcal D&\simeq\Sigma_{V,\omega}^*
\end{aligned}
\right.
\tag{21}
\end{equation}
où $\mathcal D$ est le demi-plan de Poincaré, i.e. \deleted{l'ens} la variété complexe des $z\in\mathbb C$ tels que $\deleted{\unsure{\Im z}}\ \Im z>0$
\begin{equation}
\mathcal D=\{z\in\mathbb C\mid\Im z>0\},
\tag{22}
\end{equation}
identifié à $\Sigma_{\mathbb R^2,\omega_0}^*$ en identifiant $z\in\mathcal D$ à la droite
\begin{equation}
L_z\deleted{\unsure{R_z}}=\mathbb C\cdot\underbrace{(z,1)}_{ze_1+e_2}.
\tag{23}
\end{equation}

\margin{$(e_1,e_2)$ base canonique de $\mathbb R^2$. Si on associait à $z\in\mathbb C$ la droite $\mathbb C(1,z)$, alors $\Sigma_{\mathbb R^2,\omega_0}^*$ correspondrait au demi-plan $\mathcal D^-$ au lieu de $\mathcal D=\mathcal D^+$.}

\margin{Ainsi, les classes d'isomorphie de courbes elliptiques avec rigidif. de Torelli sont en corr. biuniv. \unsure{avec l'ens} des pts de $\mathcal D$.}

Sur l'ens des classes d'isomorphie de courbes elliptiques $E=(\unsure{\Pi|H},\omega,L\subset H\otimes_{\mathbb Z}\mathbb C)$, avec rigidification de Torelli $\mathbb Z^2\simeq H$, le groupe $\mathrm{Sl}(2,\mathbb Z)$ opère à droite de façon évidente, par composition \inserted{à droite} de (18) avec les automorphismes de $\mathbb Z^2$. \deleted{\unsure{Comme on}}
